% Symbol macros for the network diagrams.
%
% Convention: lowercase for variables and functions; uppercase is reserved for
% matrices, except for integer counts, which are always N_{quantity}.
%
% This is the ONE place a symbol is renamed -- the figure, the notation table and
% any caption all pick the change up.  Numbers live in dims.tex (generated).

% --- data space --------------------------------------------------------------
\newcommand{\coordx}{x}                 % data-space node coordinate
\newcommand{\Nx}{N_{x}}                 % number of data-space nodes
\newcommand{\datavec}{u}                % data features at one timestep
\newcommand{\Nu}{N_{u}}                 % number of model input variables
\newcommand{\Nfield}{N_{F}}             % encoded fields
\newcommand{\Nfc}{N_{c_F}}              % latent channels per field
\newcommand{\Npass}{N_{\mathrm{dyn}}}    % dynamic forcings
\newcommand{\Nstatic}{N_{\mathrm{stat}}} % static forcings, given once
\newcommand{\NT}{N_{T}}                 % observed input timesteps
\newcommand{\NTf}{N_{T_f}}              % forcing-only target timesteps
\newcommand{\attrx}{a_{x}}              % data-node attributes
\newcommand{\Nax}{N_{a_{x}}}
\newcommand{\inputvec}{\psi}            % assembled data-node input
\newcommand{\Npsi}{N_{\psi}}

% --- latent space ------------------------------------------------------------
\newcommand{\coordz}{z}                 % latent-space node coordinate
\newcommand{\Nz}{N_{z}}                 % number of latent-space nodes
\newcommand{\attrz}{a_{z}}              % latent-node attributes
\newcommand{\Naz}{N_{a_{z}}}
\newcommand{\latentvec}{\varphi}        % latent state vector
\newcommand{\Nphi}{N_{\varphi}}         % number of latent channels

% --- graph edges -------------------------------------------------------------
% The encoder and decoder graphs are distinct edge sets with distinct attributes
% and distinct weights; only the widths happen to coincide, so one symbol serves
% both and the edge counts carry the distinction.
\newcommand{\attre}{a_{e}}              % edge attributes
\newcommand{\Nae}{N_{a_{e}}}
\newcommand{\NEE}{N_{EE}}               % encoder edges, data -> latent
\newcommand{\NDE}{N_{DE}}               % decoder edges, latent -> data
\newcommand{\NPE}{N_{PE}}               % processor edges, latent -> latent
\newcommand{\nbhd}{\mathcal{N}}         % incoming neighbourhood of a destination node

% --- model output ------------------------------------------------------------
\newcommand{\outvec}{y}                 % model output in data space
\newcommand{\Nout}{N_{y}}               % number of output channels
\newcommand{\idxprog}{\mathcal{P}}      % prognostic output channels
\newcommand{\idxdiag}{\mathcal{D}}      % diagnostic output channels
% Counts of the two, so the sets keep the set symbols and the integer counts are
% written N_{quantity} like every other count here.  Subscripting the set macro
% means renaming a set renames its count with it.
\newcommand{\Nprog}{N_{\idxprog}}
\newcommand{\Ndiag}{N_{\idxdiag}}
\newcommand{\Nnet}{N_{y_{\mathrm{net}}}}   % output channels emitted by the decoder
\newcommand{\Nphys}{N_{y_{\mathrm{phys}}}} % output channels diagnosed by physics

% --- processor ---------------------------------------------------------------
\newcommand{\NL}{N_{L}}                 % processor layers

% --- attention ---------------------------------------------------------------
\newcommand{\NH}{N_{H}}                 % attention heads
\newcommand{\Nc}{N_{c}}                 % channels per head, N_phi / N_H
\newcommand{\query}{q}
\newcommand{\keyv}{k}
\newcommand{\valuev}{v}
\newcommand{\edgeemb}{\epsilon}         % embedded edge attributes
\newcommand{\selfv}{r}                  % lin_self branch
\newcommand{\attn}{\alpha}

% --- mlp ---------------------------------------------------------------------
% anemoi's own name for this block (cf. the mlp_hidden_ratio config knob)
\newcommand{\ff}{f}                     % the MLP
\newcommand{\Nf}{N_{f}}                 % MLP hidden width

% --- shape annotation --------------------------------------------------------
% One macro so the styling of every wire label changes in one place.
\newcommand{\shape}[1]{\ensuremath{#1}}
\newcommand{\shapetwo}[2]{\shape{#1 \!\times\! #2}}
\newcommand{\shapethree}[3]{\shape{#1 \!\times\! #2 \!\times\! #3}}
